Este proyecto se enfoca en el desarrollo, entrenamiento y validación de un modelo inverso basado en redes neuronales profundas para la localización de fuentes neuronales a partir de señales EEG no invasivas. Se busca implementar la simulación de señales realistas mediante distintos modelos de cabeza (homogéneo, BEM, FEM) y la generación de actividad cerebral sintética con modelos biológicos como los Neural Mass Models.
    
Se entrenará un modelo de aprendizaje profundo (CNN y/o RNN) con estos datos, buscando estimar tanto la posición como la dinámica de las fuentes internas. El desempeño del modelo será evaluado cuantitativamente mediante métricas como el error de localización (LE), error temporal (TE), nMSE, PSNR, precisión, sensibilidad y F1-score. Además, se realizará una comparación sistemática con métodos clásicos como MNE, sLORETA y VARETA. 
    
También se contempla una validación cualitativa utilizando datos reales con base en evidencia clínica disponible y modalidades complementarias como fMRI o iEEG. No se incluyen en este proyecto la adquisición directa de datos en sujetos humanos ni procedimientos invasivos, enfocándose exclusivamente en simulaciones y el uso de datos estandarizados.
